\section{A nonintervening Gaussian sensing law}
\label{sec:gaussian}
In this section the task is not randomized by a physical intervention. At a productive call a current hidden sign $E\in\{-1,1\}$ emits
\begin{equation}\label{eq:sensor}
 X=aE+\sigma G,\qquad a,\sigma>0,
\end{equation}
where $G$ is a fresh standard normal variable. The audit is that current sign $E$. All history-dependent hidden sign evolutions are allowed, including the static signs and an independently refreshed fair sign. No command may alter the sign between the sensor call and its audit. The heralding interface is that of Lemma~\ref{lem:gate}; all three indicated sign models with constant gate rate $p_0$ are admitted. Thus arbitrary old observations or gate histories may be informative, but the class has no positive uniform conditional variance floor.

Under the independent fair reference, the actual productive density and conditional moments are
\begin{equation}\label{eq:gaussian-moments}
 \nu(dx)=\frac{\varphi_\sigma(x-a)+\varphi_\sigma(x+a)}2dx,
 \quad f(x)=\tanh(kx),\quad w(x)=\operatorname{sech}^2(kx),
 \quad k=a/\sigma^2.
\end{equation}
The formulas follow by dividing the two Gaussian likelihoods; their ratio is $\exp(2kx)$. For $b\ge1$, let $t_b\ge0$ be the unique threshold with $\nu(|X|\ge t_b)=1/b$, taking $t_1=0$, and define
\begin{equation}\label{eq:D}
 D(b)=b\int_{|x|\ge t_b}w(x)\nu(dx),\qquad
 U(b)=2b\int_{|x|\ge t_b}(1-|f(x)|)\nu(dx).
\end{equation}

The two-sided threshold construction and the logarithmic waiting scale were already proved for a positive-prior-floor class in \cite{GTF52}. Here the lower is instead supplied by one admitted fair reference member: this removes that floor, permits static hidden signs, and yields a ratio-one minimax asymptotic uniformly over every larger terminal alphabet. The conditional-variance selection principle below is standard in regression with rejection \cite{DHZ}, with the earlier recognition tradeoff in \cite{Chow}. The assertion here additionally constrains the paid stopped experiment and all of its persistent states; it does not assert priority for a threshold rule.

\begin{theorem}[Terminal-value transition and the sharp waiting asymptotic]\label{thm:gaussian}
Put $b=N/\kappa$, with $\kappa=1+1/p_0$, and assume $N\ge\kappa$. One terminal value gives worst risk one. In particular,
\begin{equation}\label{eq:gaussian-width}
 R_\st^\Theta(N,3)=R_\cp^\Theta(N,1)=1.
\end{equation}
For each $R$ in the collection
\[
 \{R_\st^\Theta(N,M):M\ge4\}\ \cup\
 \{R_\evt^\Theta(N,M):M\ge2\}\ \cup\
 \{R_\cp^\Theta(N,m):m\ge2\},
\]
one has $D(b)\le R\le U(b)\le2D(b)$, and uniformly over this entire collection,
\begin{equation}\label{eq:gaussian-asymp}
 R\sim 4\exp(-2a^2/\sigma^2)
       \exp\!\left[-\frac{2a}{\sigma}\sqrt{2\log(N/\kappa)}\right]
 \quad\text{as }N/\kappa\longrightarrow\infty.
\end{equation}
The state-stopped class is empty below three labels; the event-stopped class is empty below two. Two event-stopped labels already permit both live ports and both terminal values. The asymptotic is not a fixed-budget limit as $p_0\downarrow0$.
\end{theorem}
\begin{proof}
The lower uses the independently refreshed fair sign, not a chosen posterior prior for each competitor. Lemma~\ref{lem:lower} bounds the terminal $X$-law by $b\nu$, and the conditional residual variance in \eqref{eq:gaussian-moments} is unavoidable even with infinitely many output values. Since $w$ decreases with $|x|$, its minimum integral against a mass-one density bounded by $b$ is $D(b)$. For completeness, subtract the threshold rule from any feasible density: on $|x|\ge t_b$, $w(x)\le w(t_b)$ and $h-b\le0$; on the complement, $w(x)\ge w(t_b)$ and $h\ge0$. Integrating the product and using zero total mass difference proves this minimization without a restriction to threshold competitors.

For the upper, accept exactly when $|X|\ge t_b$ and output $\operatorname{sign}(X)$. Reflection of the Gaussian kernel shows that its acceptance probability is $1/b$ for each fixed current sign, and its unnormalized squared-loss contribution is $U(b)/b$ for each sign. The same identities hold for every mixture of signs conditional on any earlier history. The proof of Lemma~\ref{lem:upper} therefore gives expected cost at most $\kappa b$ and risk $U(b)$, without refreshing or estimating the hidden sign. This uses two live port roles and the two outputs $\pm1$, packed into four state-stopped or two event-stopped labels by Lemma~\ref{lem:packing}. Algebra gives
\[
 \frac{2(1-|f|)}{1-f^2}=\frac2{1+|f|},
\]
so $D(b)\le U(b)\le2D(b)$ and, when $t_b>0$,
$U(b)/D(b)\le2/(1+\tanh(kt_b))$.

If there is only one terminal value $r\in[-1,1]$, the two admitted static-sign models give worst-case risk at least
$\max\{(1-r)^2,(-1-r)^2\}\ge1$, irrespective of stopping. Outputting zero at the first productive report attains one. The mandatory physical ports leave only one terminal value for a state-stopped observer when $M=3$. This subtraction is not valid for the event-stopped observer, whose two labels can carry the two different roles on different branches. This argument allows all fresh randomized stopping and routing, not only thresholds.

It remains to identify the asymptotic constant. Let $u_b=(t_b-a)/\sigma$ and $c=a/\sigma>0$. The tail equation is
\[
 b^{-1}=\overline\Phi(u_b)+\overline\Phi(u_b+2c).
\]
Integration by parts gives, for $u>0$,
$\frac{u}{1+u^2}\varphi(u)\le\overline\Phi(u)\le\varphi(u)/u$.
The second tail is negligible relative to the first, so
$u_b=\sqrt{2\log b}+O(\log\log b/\sqrt{\log b})$.
In particular $t_b=a+\sigma\sqrt{2\log b}+o(1)$. For any fixed $v>0$, the same Gaussian tail bounds show
$\nu(|X|>t_b+v)/\nu(|X|\ge t_b)\to0$. Thus $|X|-t_b\to0$ in probability under the selected tail law. Moreover
$0\le w(X)/w(t_b)\le1$ there, and this ratio tends to one in that probability. Bounded convergence in probability yields $D(b)/w(t_b)\to1$. Since $w(t_b)\sim4e^{-2kt_b}$, the right side of \eqref{eq:gaussian-asymp} follows. Finally $U(b)/D(b)\to1$. The lower and upper do not depend on the available dictionary width once it is at least two, proving the claimed uniformity under the three listed budget conventions.
\end{proof}

This application is also a raw transformed-task verification of Theorem~\ref{thm:main}. After a sensor report, choose a fresh fair $g\in\{-1,1\}$, use virtual pair $(gX,gE)$ and export $T_gr=gr$. Averaging yields exactly the reference joint law for every pre-sensor sign distribution, with $D_0=2$. No physical randomization of the target occurs. The symmetric threshold rule above removes even the virtual coin. Its sharp asymptotic arises from the actually acquired Gaussian tails, not a formal one-dimensional parameter space. The one-value obstruction requires the two static models in addition to the fair lower witness; it is not inferred from a variance-floor hypothesis.
